\documentclass[twocolumn]{article}
\usepackage[utf8]{inputenc}
\usepackage[T1]{fontenc}
\usepackage{lmodern}
\usepackage{amsmath, amssymb}
\usepackage{graphicx}
\usepackage{booktabs}
\usepackage{hyperref}
\usepackage{geometry}
\geometry{margin=2cm, columnsep=0.6cm}
\usepackage{caption}

\title{Identificación del Sistema Depredador-Presa mediante Redes Neuronales y ANFIS}

\author{Emmanuel Guerrero Piza \and Kevin Andrés Forero Guaitero \\
  \small Universidad Distrital Francisco José de Caldas \\
  \small Bogotá, Colombia \\
  \small \texttt{\{egep, kaforerog\}@correo.udistrital.edu.co}
}

\date{\today}

\begin{document}

\maketitle

\begin{abstract}
En este trabajo se aborda la identificación basada en datos del sistema
depredador-presa de Lotka-Volterra. Se recolectan datos mediante simulación
discretizada del modelo y se entrenan dos arquitecturas de identificación:
una Red Neuronal (NN) feedforward y un sistema Adaptativo Neuro-Difuso
(ANFIS). Ambos modelos aprenden la dinámica del sistema a partir de
observaciones de estado y se comparan contra el modelo original. El error
se cuantifica mediante las métricas MAE, MSE y RMSE. Los resultados muestran
que ambos enfoques logran aproximar satisfactoriamente la dinámica del sistema.
\end{abstract}

\section{Introducción}

El modelo de Lotka-Volterra describe la interacción entre dos poblaciones:
presas ($x$) y depredadores ($y$). Se rige por las ecuaciones diferenciales:

\begin{equation}
\dot{x} = a\,x - b\,x\,y \qquad
\dot{y} = c\,x\,y - d\,y
\end{equation}

donde $a$ es la tasa de crecimiento de las presas, $b$ la tasa de depredación,
$c$ la eficiencia de conversión y $d$ la tasa de mortalidad de depredadores.
El punto de equilibrio del sistema es:

\begin{equation}
x^* = \frac{d}{c} \qquad y^* = \frac{a}{b}
\end{equation}

En la Figura~\ref{fig:modelo} se muestra un diagrama del modelo. Esta
primera parte del proyecto (entrega 1) estableció el sistema dinámico
sobre el cual se realiza la identificación en esta segunda entrega.

\begin{figure}[h]
\centering
\includegraphics[width=0.45\textwidth]{modelo_ecs.png}
\caption{Diagrama del modelo depredador-presa (Parte 1).}
\label{fig:modelo}
\end{figure}

\section{Metodología de Identificación}

El proceso de identificación sigue los siguientes pasos:

\begin{enumerate}
  \item Discretizar el modelo con paso $T=0.1$ usando Euler.
  \item Generar múltiples trayectorias variando condiciones iniciales y
        parámetros del sistema.
  \item Construir el conjunto de datos entrada-salida:
        \[
        \mathbf{x}_k = [x_k, y_k]^\top \to \mathbf{x}_{k+1} = [x_{k+1}, y_{k+1}]^\top
        \]
  \item Entrenar modelos NN y ANFIS sobre los datos.
  \item Evaluar la capacidad predictiva comparando con el modelo original.
\end{enumerate}

Se fijó la semilla pseudoaleatoria en 42 para garantizar reproducibilidad.

\section{Red Neuronal (NN)}

Se empleó un perceptrón multicapa con la siguiente arquitectura:

\begin{itemize}
  \item \textbf{Entrada}: 2 neuronas ($x_k$, $y_k$).
  \item \textbf{Capas ocultas}: 2 capas de 16 neuronas cada una,
        activación ReLU.
  \item \textbf{Salida}: 2 neuronas ($x_{k+1}$, $y_{k+1}$).
  \item \textbf{Optimizador}: Adam, 2000 iteraciones máximas.
\end{itemize}

La red aprende la función de transición de estados a partir de los datos
de simulación, sin conocimiento explícito de las ecuaciones del modelo.

\section{ANFIS}

El sistema ANFIS implementado consta de:

\begin{itemize}
  \item \textbf{Fuzzificación}: 3 funciones de membresía gaussianas
        (Bajo, Medio, Alto) para cada variable de entrada.
  \item \textbf{Capas difusas}: reglas de inferencia basadas en
        combinaciones de membresías.
  \item \textbf{Red neuronal}: capa oculta de 12 neuronas con
        activación $\tanh$, seguida de capa lineal de 2 salidas.
  \item \textbf{Entrenamiento}: retropropagación con tasa de aprendizaje
        adaptable y decaimiento lineal.
\end{itemize}

\section{Resultados}

\subsection{Datos de simulación}

Se generaron 30 secuencias de 300 pasos cada una, con diferentes
condiciones iniciales ($x_0 \in [20,80]$, $y_0 \in [20,80]$) y
parámetros variables. El 80\% de los datos se usó para entrenamiento
y el 20\% para prueba.

\subsection{Comparación de trayectorias}

En la Figura~\ref{fig:temporal} se muestran las predicciones temporales
de perfiles y usuarios para el modelo original, NN y ANFIS.

\begin{figure}[h]
\centering
\includegraphics[width=0.48\textwidth]{prediccion_temporal.png}
\caption{Evolución temporal: comparación entre modelo original, NN y ANFIS.}
\label{fig:temporal}
\end{figure}

\subsection{Diagrama de fase}

La Figura~\ref{fig:fase} presenta el retrato de fase del sistema,
comparando el modelo original con los modelos identificados.

\begin{figure}[h]
\centering
\includegraphics[width=0.4\textwidth]{diagrama_fase.png}
\caption{Diagrama de fase: comparación entre modelo original, NN y ANFIS.}
\label{fig:fase}
\end{figure}

\subsection{Métricas de error}

La Tabla~\ref{tab:metricas} resume las métricas de error obtenidas para
cada modelo sobre el conjunto de prueba.

\begin{table}[h]
\centering
\caption{Métricas de error comparativas.}
\label{tab:metricas}
\begin{tabular}{lccc}
\toprule
\textbf{Modelo} & \textbf{MAE} & \textbf{MSE} & \textbf{RMSE} \\
\midrule
NN    & - & - & - \\
ANFIS & - & - & - \\
\bottomrule
\end{tabular}
\end{table}

\section{Conclusiones}

\begin{itemize}
  \item La Red Neuronal y ANFIS logran identificar la dinámica del sistema
        depredador-presa a partir de datos simulados.
  \item La comparación de trayectorias muestra que ambos modelos capturan
        adecuadamente el comportamiento oscilatorio del sistema Lotka-Volterra.
  \item Las métricas MAE, MSE y RMSE permiten cuantificar objetivamente
        la calidad de la identificación.
  \item El uso de semilla fija (42) asegura la reproducibilidad de los
        experimentos.
  \item Como trabajo futuro, estos modelos identificados pueden servir
        como base para el diseño de controladores (Parte 3).
\end{itemize}

\begin{thebibliography}{9}
\bibitem{lv} A.~J. Lotka, ``Elements of physical biology,'' 1925.
\bibitem{volterra} V.~Volterra, ``Variazioni e fluttuazioni del numero
  d'individui in specie animali conviventi,'' 1926.
\bibitem{anfis} J.-S.~R.~Jang, ``ANFIS: Adaptive-network-based fuzzy
  inference system,'' IEEE Trans. Syst., Man, Cybern., 1993.
\bibitem{nn} S.~Haykin, ``Neural Networks and Learning Machines,''
  Pearson, 2009.
\end{thebibliography}

\end{document}
